\section{Introduction}

Since the 1998 discovery \citep{fukuda} of neutrino oscillations,
there have been major advances in neutrino disappearance
and appearance  oscillation measurements
% of solar neutrinos\citep{SNO},
%accelerator based $\nu_{\mu}$ disappearance\citep{K2K,MINOS} and
%$\nu_{\mu}\rightarrow\nu_{e}$ appearance\citep{LBNE-theta13} and
%reactor based $\nu_{e}$ disappearance\citep{Reactor-theta13}. These
%measurements largely determine 
and all the fundamental neutrino mixing parameters \citep{PDG-2016} have been determined
% in the Pontecorvo-Maki-Nakagawa-Sakata (PMNS) mixing matrix\citep{PDG-2016}
except for the mass hierarchy and the charge-parity (CP)
phase $\delta_{CP}$. 
Evidence of $\delta_{CP}\ne 0, \pi$  leads to the non-conservation or violation
of the charge-parity symmetry (CPV).
This is tested by measuring the neutrino $\nu_{\mu}\rightarrow\nu_{e}$ and antineutrino
$\overline{\nu}_{\mu}\rightarrow\overline{\nu}_{e}$ appearance oscillation
event rates to determine if the neutrino and antineutrino oscillation appearance probabilities, 
$P\left(\nu_{\mu}\rightarrow\nu_{e}\right)$
and
$\bar{P}\left(\overline{\nu}_{\mu}\rightarrow\overline{\nu}_{e}\right)$
are equal in vacuum \citep{msw}
at the same ratio of  
the oscillation distance $L$
over
the neutrino energy $E$ or $\frac{L}{E}$.
%If the measured probabilities are unequal in vacuum\citep{msw} at the same
%value of $E/L$ where $E$ is the neutrino/antineutrino energy, then CPV is established.
Major %experiental efforts 
long-baseline neutrino experiments \citep{t2k+nova} have been built
and future projects \citep{dune+hyperk} are proposed to determine
these probabilities  using separate $\nu_{\mu}$ and $\overline{\nu}_{\mu}$ beams
that cross near and far detectors. 
The probabilities are obtained from near detector measurements of the 
%unoscillated  
$\nu_{\mu}+N$ and $\overline{\nu}_{\mu}+N$ charged current (CC) interactions and cross sections,
where $N$ is the target nucleon, 
and far detector measurements of $\nu_{e}+N$ and $\overline{\nu}_{e}+N$ CC interactions. 
%The neutrino and antineutrino oscillation appearance probabilities
%can be determined from efficiency corrected ratios of numbers of events 
%of $n(\nu_e)/n(\nu_\mu)$ and $n(\bar{\nu}_e)/n(\bar{\nu}_\mu)$,
%respectively. These ratios would be equal if CP is conserved in neutrino mixing.

In this paper, the T2K Collaboration, using the off-axis near detector (ND280), presents
a measurement at a peak energy %\textasciitilde{}
$\sim$0.6 GeV of the %unoscillated
charged current inclusive (CCINC)
$\nu_{\mu}+N$ cross section and first CCINC
measurements of the $\overline{\nu}_{\mu}+N$ cross section and
their ratio of the  $\overline{\nu}_{\mu}+N$ over the $\nu_{\mu}+N$
CCINC cross section.
These $\nu_{\mu}$  and $\bar\nu_{\mu}$ measurements 
%and the precision on these measurements in particular
are important to understand their
impact on future CPV measurements and to 
test neutrino cross section models.

%Besides measuring the electron or antielectron neutrino appearance
%events at the far detector,  the unoscillated $\nu_\mu$ and
%$\bar{\nu}_\mu$ rates must be measured to make complete 
%measurement of the neutrino and antineutrino appearance
%probabilities.  
%A necessary component of the unoscillated rates are near detector measurements of 
%the neutrino and antineutrino charged current inclusive (CCinc) and
%charged current quasi-elastic (CCQE)
%cross sections, the interaction rates and their ratios.
%In this paper we discuss measurements of the CCinc 
%$\nu_\mu$ and $\bar{\nu}_\mu$ cross sections and their ratio. 
T2K has published flux averaged neutrino-mode  measurements  of CCINC \citep{t2k-carbon-ccinc}
and charged current quasi-elastic like (CCQE) \citep{t2k-cceq}
cross sections per nucleon of  
 $(6.91\pm0.13{\rm (stat.)}\pm0.84{\rm (syst.)}) \times 10^{-39}$cm$^2$ 
and $(4.15\pm0.6) \times 10^{-39}$cm$^2$, 
respectively. These measurements were performed using the
Fine-Grain Detector (FGD) which has different detector systematics
compared to the measurements presented in this paper.
There are no published CCINC $\bar\nu_{\mu}$
measurements at
energies below 1.5 GeV, however
the MINVERVA Collaboration recent published \citep{minerva+nu+antinu}
CCINC results above 2 GeV
and
the MiniBooNE Collaboration has published \citep{miniboone} CCQE
measurements in both 
%$\nu_\mu$  and $\bar\nu_\mu$  
\nubarandnu
modes which require
larger axial mass values compared to other experiments to fit their observed data. There are several
multinucleon models (2 particle 2 hole, or 2p2h)
 \citep{martini-1,nieves-1,martini-2} proposed to explain large cross sections.
In addition, in some models it has been predicted \citep{martini-1} that the difference between the 
%neutrino and antineutrino
\nuandnubar 
cross sections is expected to increase when 
%multinucleon interaction 2 particle and 2 hole (2p2h) nuclear model 
2p2h effects \citep{2p2h-refs}
are included. The measurements of the ratio, sum, and difference of these cross sections,  which
have very different systematic errors, will be presented.

Following this introduction, the paper is organized as follows. We
begin with a description of the ND280 off-axis detector and the neutrino
beam in Section II. Then the Monte Carlo (MC) simulation is presented in
Section III, followed by the event selection given in Section IV.
The analysis methods and systematic error evaluations are presented
in Sections V and VI and we finally conclude with the Results and
Conclusions in Sections VII and VIII.


